\subsection{伽罗瓦扩张的定义}

定理 1.——设 N 是 K 的一个代数扩张，\(\Gamma\) 是 N 的 K-自同构的群。则下列断言等价：

\begin{enumerate}
\def\labelenumi{\alph{enumi})}
\item
  \(\mathbf{N}\) 中在 \(\Gamma\) 下不变的每个元素都属于 \(\mathbf{K}\) 在 \(\mathbf{N}\) 中的像。
\item
  N 是 K 的一个可分拟伽罗瓦扩张。
\item
  对每个 \(x \in \mathbf{N}\) ，\(x\) 在 \(\mathbf{K}\) 上的极小多项式在 \(\mathbf{N}[\mathbf{X}]\) 中分解为一些互异的 1 次多项式之积。
\end{enumerate}

b) 与 c) 的等价性由命题 6 的推论（V，第 40 页）及拟伽罗瓦扩张的定义（V，第 53 页，定义 2）得出。我们把 K 与它在 N 中的典范像等同起来。

\(a) \Rightarrow c)\) ：设 K 是 \(\Gamma\) 的不变量域。设 \(x \in \mathbf{N}\) ，其在 K 上的极小多项式为 \(f\) ，并设 A 为 \(f\) 在 N 中的全体根组成的集。令

\[
g (\mathbf {X}) = \prod_ {y \in \mathsf {A}} (\mathbf {X} - y).
\]

每个自同构 \(\sigma \in \Gamma\) 都诱导 \(\mathbf{A}\) 的一个置换，从而保持多项式 \(g \in \mathbf{N}[\mathbf{X}]\) 的系数不变。于是我们有 \(g \in \mathbf{K}[\mathbf{X}]\) ，且由于 \(g(x) = 0\) ，多项式 \(g\) 是 \(f\) 在 \(\mathbf{K}[\mathbf{X}]\) 中的倍式（V，第 16 页，定理 1）。再者，\(f\) 与 \(g\) 都是首一的，且 \(g\) 整除 \(f\) （IV，第 16 页，命题 5）；从而 \(f = g\) ，这意味着极小多项式 \(f\) （\(x\) 的，在 \(\mathbf{K}\) 上的）在 \(\mathbf{N}[\mathbf{X}]\) 中是一些互异的 1 次多项式之积。

\(c) \Rightarrow a)\) ：设 \(x\) 是 \(\mathbf{N}\) 中不属于 \(\mathbf{K}\) 的元素。用 \(\Omega\) 表示 \(\mathbf{K}\) 的一个把 \(\mathbf{N}\) 作为子扩张包含在内的代数闭包（V，第 23 页，定理 2）。设 \(f\) 是 \(x\) 在 \(\mathbf{K}\) 上的极小多项式，由假设它的次数 \(\geqslant 2\) ，并设 \(\mathbf{A}\) 是 \(f(\mathbf{X})\) 在 \(\mathbf{N}\) 中的根集。若条件 \(c)\) 成立，我们有 \(f(\mathbf{X}) = \prod_{y \in \mathbf{A}} (\mathbf{X} - y)\) ，于是

（V，第 53 页，推论 1）A 是 \(x\) 在 \(\Omega\) 中的共轭元之集。由于 \(f\) 的次数 \(\geqslant 2\) ，A 中存在元素 \(y \neq x\) ，从而存在一个 K-自同构 \(u\) （\(\Omega\) 的）使得 \(u(x) = y\) 。而在假设 \(c\) ) 下，K 的扩张 N 是拟伽罗瓦的，故 \(u(N) = N\) （V，第 54 页，推论 1）；由此 \(u\) 诱导 N 的一个 K-自同构 \(\sigma\) ，使得 \(\sigma(x) = y \neq x\) ，因此 K 是 \(\Gamma\) 的不变量域。

定义 1.——K 的扩张 N 称为伽罗瓦扩张，如果它是代数的且满足定理 1 的等价条件 a)、b)、c)。

设 N 是一个域，\(\Gamma\) 是 N 的一个自同构群，\(N_{0}\) 是 \(\Gamma\) 的不变量域。当 N 在 \(N_{0}\) 上代数时，它是 \(N_{0}\) 的伽罗瓦扩张。但情形未必总是如此：例如，设 K 是无限域，取 N 为有理分式域 \(\mathrm{K}(\mathrm{X})\) ；对每个 \(a \in K\) ，令 \(\sigma_{a}\) 为 \(\mathrm{K}(\mathrm{X})\) 的把 \(f(\mathrm{X})\) 映到 \(f(\mathrm{X} + a)\) 的自同构。所有 \(\sigma_{a}\) 组成的集是 \(\mathrm{K}(\mathrm{X})\) 的一个自同构群，其不变量域容易看出是 K；然而 \(\mathrm{K}(\mathrm{X})\) 在 K 上不是代数的。

设 \(\Omega\) 是 K 的一个代数闭包，设 A 表示由 \(\Omega\) 中在 K 上可分的元素组成的集，B 表示 A 的元素在 K 上的共轭元组成的集。则 B 由在 K 上代数且可分的元素组成。因此（V，第 39 页，命题 6 及第 56 页，命题 5）域 K(B) 是 K 的可分拟伽罗瓦扩张；换句话说，由 A 生成的拟伽罗瓦扩张（V，第 55 页）是 K 的伽罗瓦扩张；我们也称它为由 \(\Omega\) 的子集 A 生成的 K 的伽罗瓦扩张。

特别地，\(\Omega\) 中一族在 K 上可分的多项式的分裂域，以及 K 的可分闭包，都是 K 的伽罗瓦扩张。

命题 1.——设 N 是 K 的一个扩张，\((\mathbf{N}_{i})_{i\in I}\) 是 N 的子扩张的一个非空族。令 \(E=\bigcap_{i\in I}N_{i}\) 且 \(F=K\left(\bigcup_{i\in I}N_{i}\right)\) 。若所有扩张

\(N_{i}\) 在 K 上都是伽罗瓦的，则 E 与 F 亦然。

首先，E 在 K 上是代数且可分的（V，第 36 页，命题 1），对 F 同样成立（V，第 41 页，命题 8）。此外，由命题 4（V，第 55 页），E 与 F 在 K 上都是拟伽罗瓦的。

命题 2.——设 N 是 K 的一个伽罗瓦扩张，E 是 N 的一个子扩张，且在 K 上次数有限。存在 N 的一个包含 E 的子扩张 F，它在 K 上是伽罗瓦的且次数有限。

由于 N 在 K 上拟伽罗瓦，V 第 56 页推论 1 证明了存在 N 的一个包含 E 且在 K 上次数有限的拟伽罗瓦子扩张 F。由于 N 在 K 上可分，对 F 同样成立（V，第 36 页，命题 1），故 F 在 K 上是伽罗瓦的。

命题 2 蕴含下述结果：设 \(\Omega\) 是 K 的一个代数闭包，\(E_1, \ldots, E_n\) 是包含在 \(\Omega\) 中的 K 的有限次可分代数扩张。则存在 K 的一个有限次伽罗瓦扩张 N，它包含在 \(\Omega\) 中且包含 \(E_1, \ldots, E_n\) .

